\subsection*{The complete $3$-adic common-image obstruction}

\statement{B.64}{Proved}{An Eichler centralizer is fully realized in the common $3$-adic image}
Let $q_\infty:\widehat P\to\operatorname{PSp}_{10}(\mathbb Z_3)$
be the branched-cover homology action.
Let $N_h$ be the normal closure of B.56's $h$ in $H_T$,
put $R_T=N_h\cap[H_T,H_T]$, and use closures in this
$3$-adic matrix target in the following assertion.
Let $\mathcal I$ be the identity fibre of reduction
modulo three in $\operatorname{PSp}_{10}(\mathbb Z_3)$.
Then
\[
 \overline{q_\infty(R_T)}\cap
 \overline{q_\infty([H_5,H_5])}\cap\mathcal I
       =C_{\operatorname{PSp}_{10}(\mathbb Z_3)}(A_5)
                      \cap\mathcal I
       =:\mathcal U.                                      \tag{B.64.1}
\]
The group $\mathcal U$ has the explicit description
\[
 \mathcal U\cong
 \left\{
 \begin{pmatrix}1+3a&3b\\9c&1+3d\end{pmatrix}:
 a,b,c,d\in\mathbb Z_3,\quad
 (1+3a)(1+3d)-27bc=1
 \right\}.                                                \tag{B.64.2}
\]
It is a nontrivial three-dimensional $3$-adic group.
The actual residual group has image contained in
$\mathcal U$ by B.62, but (B.64.1) does not give a
nonidentity element of $\mathcal K$.
It determines a common image in one representation,
with compatible lifts to its two source completions;
equality of those lifts in $\widehat P$ remains unproved.

\begin{proof}
\emph{The actual integral $A_5$ lattice.}
In the twelve-strand unreduced Burau module at $-1$,
let $\epsilon=(1,-1,\ldots,1,-1)$ and
$\mathbf1=(1,\ldots,1)^t$.
The closed branched-cover homology lattice is
\[
              \Lambda=\ker\epsilon/\mathbb Z\mathbf1.
                                                               \tag{B.64.3}
\]
This is the even-strand version of B.62's branched-cover
chain construction: the disk cover has two boundary
components, and capping kills its primitive boundary
class.  Representatives with last coordinate zero
give the basis with inclusion columns
$(e_j,-\epsilon_j,0)^t$, for $1\leq j\leq10$.
The action of a twelve-dimensional matrix $B$ on this
basis is obtained by subtracting its last output
coordinate from the first ten.
All eleven Artin generators preserve the resulting
integral unimodular symplectic form.

Apply this construction to the actual $\beta_2,\beta_3$
certified in B.8.  Negate the first resulting matrix.
The two matrices then satisfy
$S^2=R^3=(SR^{-1})^5=I$ and generate an actual copy
of $A_5$, of order sixty, in $\operatorname{Sp}(\Lambda)$.
Their character on orders $1,2,3,5$ is
\[
                         (10,2,-2,0).                       \tag{B.64.4}
\]
It is twice the irreducible degree-five character.
For example that character is the augmentation
character of the six-point action; its squared norm
is $(25+15+20)/60=1$.
Thus over characteristic zero the symplectic
centralizer of $A_5$ is $\operatorname{SL}_2$.
This identification is also proved by the explicit
commuting algebra below, without relying on a modular
semisimplicity assertion.

The exact certificate
\nolinkurl{next_3adic_symplectic/icosahedral_integral_centralizer.py}
solves the commuting equations over $\mathbb Q$, then
takes their saturated lattice at three.
The commuting order has rank four and its symplectic
Lie lattice has rank three.
The certificate provides a basis $H,E,F$ of the latter
over $\mathbb Z_3$ satisfying
\[
 \begin{gathered}
 H^2=I,\qquad E^2=F^2=0,\\
 [H,E]=2E,\qquad [H,F]=-2F,\qquad [E,F]=3H,\\
 EF=\frac{3(I+H)}2,\qquad FE=\frac{3(I-H)}2.
 \end{gathered}                                             \tag{B.64.5}
\]
All entries and the basis changes are rational numbers
with denominator prime to three.
The four matrices $I,H,E,F$ are independent modulo
three in $\operatorname{Mat}_{10}(\mathbb Z_3)$.
They therefore span the full saturated commuting
order, not a suborder of index divisible by three.
The complete equations, bases, multiplication checks
and actual output are in the accompanying JSON
certificate and standard-output file.

Equations (B.64.5) identify this order with
\[
 \mathcal O=
 \begin{pmatrix}\mathbb Z_3&\mathbb Z_3\\
                 3\mathbb Z_3&\mathbb Z_3\end{pmatrix},
 \qquad
 H\longleftrightarrow\begin{pmatrix}1&0\\0&-1\end{pmatrix},
 \quad E\longleftrightarrow E_{12},
 \quad F\longleftrightarrow3E_{21}.                         \tag{B.64.6}
\]
The symplectic adjoint fixes $I$ and negates $H,E,F$,
so its norm on this order is the ordinary two-by-two
determinant.  The integral symplectic centralizer is
therefore $\mathcal O^1$, its determinant-one group.
Projective centralization introduces no additional
elements: any signs in conjugating the chosen linear
$A_5$ lift would define a homomorphism
$A_5\to\{\pm1\}$, necessarily trivial.

The natural lattice filtration must be retained.
Since $I,H,E,F$ form a primitive order basis,
congruence to $I$ in the ten-dimensional lattice
means $g-I\in3\mathcal O$.
Under (B.64.6), the lower-left entry is then divisible
by nine, not merely by three.
Thus $\mathcal O^1\cap(1+3\mathcal O)$ is precisely
(B.64.2).  Ordinary entrywise reduction of the displayed
two-by-two matrices modulo three is not faithful on
$\mathcal O/3\mathcal O$ and must not replace this
filtration.
For arbitrary $a,b,c\in\mathbb Z_3$, the determinant
condition uniquely determines
$1+3d=(1+27bc)/(1+3a)$; hence this group is homeomorphic
to $\mathbb Z_3^3$.

\emph{The two derived centralizers.}
Work in the common geometric marking of B.61.
The unlabelled $A_5$ action used above is conjugate to
it by B.52; the induced map on integral homology is
symplectic, so all lattice assertions transport.
Let $G_T$ and $G_5$ be the connected derived algebraic
groups of the symplectic centralizers of $T$ and $D$
over $\mathbb Q_3$.
The degree-five character restricts to $T\cong A_4$
as $1'\oplus1''\oplus3$.
Since homology has two copies, over an algebraic
closure its symplectic centralizer is
$\operatorname{GL}_2\times\operatorname{SL}_2$.
For $D\cong C_5$, the five characters each have
multiplicity two; its symplectic centralizer is
$\operatorname{SL}_2\times\operatorname{GL}_2
                         \times\operatorname{GL}_2$.
Consequently
\[
                    \dim G_T=6,\qquad \dim G_5=9.           \tag{B.64.7}
\]
These are dimensions of the derived groups, rather
than the full centralizers of dimensions seven and eleven.
The connected $A_5$ centralizer $\operatorname{SL}_2$
is contained in both derived groups: its image in
either abelian torus quotient is trivial.

The source lifts of $H_T,H_5$ centralize the respective
linear finite-group lifts.
Indeed any possible signs define characters
$A_4\to C_2$ or $C_5\to C_2$, both trivial.
One may use the disk braid lifts with zero linking on
the fixed pair $(4,6)$, as in B.62.
Thus every source commutator maps into the appropriate
connected derived algebraic group.
Every $n_j$ in B.63 is a product of a conjugate of $h$
and the inverse of another conjugate, so it lies in
$R_T$; every $c_j$ lies in $[H_5,H_5]$.

\emph{A matrix logarithm argument with the full lattice.}
For $G=G_T,G_5$, set
\[
 L_G=\operatorname{Lie}(G)\cap
                         \operatorname{Mat}_{10}(\mathbb Z_3),
 \qquad
 U_G=G(\mathbb Q_3)\cap
                         (I+3\operatorname{Mat}_{10}(\mathbb Z_3)).
                                                               \tag{B.64.8}
\]
The lattice $L_G$ is saturated in the matrix lattice
and has rank $\dim G$.
Convergent matrix logarithm and exponential give
\[
                U_G=\exp(3L_G),\qquad
                U_G\cap(I+3^r\operatorname{Mat}_{10})
                         =\exp(3^rL_G)\quad(r\geq1).        \tag{B.64.9}
\]
The analytic series and the powerful Lie-lattice
construction are recorded in
\citet[Proposition 6.22 and Theorem 9.8]{DE_DixonEtAl1999}.
Here $3L_G$ is powerful because
$[3L_G,3L_G]\subseteq9L_G$.

The algebraic-group comparison in (B.64.9) is direct.
Logarithm and exponential preserve the centralizer
equations.  They also identify the symplectic group
condition with its Lie equation.
Over a splitting field, the derived-group condition
on each multiplicity factor is determinant one.
Its logarithm is the trace-zero condition.
If a matrix is $I+3X$ with $X$ integral, the eigenvalues
of every restricted block of $X$ are algebraic integers.
Its determinant character therefore belongs to
$1+3\mathcal O_K$ in the splitting field, even if the
eigenspace projector has a denominator divisible by three.
The scalar logarithm is injective on this subgroup,
so determinant one and trace zero are equivalent here.
Saturation then gives exactly $3^rL_G$, including at
the ramified prime three.

Square the actual words $n_j,c_j$ before using their
linear matrices.  This removes the possible sign of
a lift whose projective mod-three image is the identity.
They belong to $U_{G_T}$ and $U_{G_5}$, respectively.
By B.63, their logarithms divided by three reduce
modulo three to the six or nine vectors $2X_j$.
These vectors are independent in the ambient matrix
lattice.  By (B.64.7) and saturation they form bases
of $L_{G_T}/3L_{G_T}$ and $L_{G_5}/3L_{G_5}$.

Here is the generation argument explicitly.
For any $L=L_G$, put $U_r=\exp(3^rL)$.
Then $U_r/U_{r+1}\cong L/3L$, and cubing sends
$\exp(3^rX)$ to $\exp(3^{r+1}X)$.
If finitely many elements of $U_1$ give a basis of
$U_1/U_2$, their $3^{r-1}$ powers give a basis of
$U_r/U_{r+1}$ for every $r$.
Successively matching an arbitrary element modulo
$U_2,U_3,\ldots$ produces a convergent product of
these powers.  Their closed generated subgroup is
therefore all of $U_1$.
Applied to the squared actual words, this proves
\[
 U_{G_T}\subseteq\overline{q_\infty(R_T)},
 \qquad
 U_{G_5}\subseteq\overline{q_\infty([H_5,H_5])},            \tag{B.64.10}
\]
using the injective projectivization on the principal
identity fibre.

The group $\mathcal U$ lies in both groups on the left
of (B.64.10), since its connected algebraic group is
contained in both derived centralizers.
Conversely, an element in the intersection of the
two source-image closures centralizes $T$ and $D$,
and therefore centralizes $\langle T,D\rangle=A_5$.
In the principal identity fibre it belongs to
$\mathcal U$.  This proves (B.64.1).

Finally, the closure of $R_T$ in $\widehat{H_T}$ maps
onto its compact $3$-adic image closure, and the same
holds for $[H_5,H_5]$ in $\widehat{H_5}$.
Thus the common image is realized compatibly at all
$3$-power levels in each source separately.
The tetrahedral source can remain in the closed
normal rigid-test kernel, and both sources in their
closed commutator subgroups.
The map from $\widehat P$ to this matrix group can
identify distinct elements, so none of these
statements implies a common nonidentity lift in
$\widehat P$ or in $\mathcal K$.
\end{proof}
